% ============================================================
%  INF4/1 — Algorithmen und Datenstrukturen
%  LEARNINGS-ZUSAMMENFASSUNG (Klausurvorbereitung)
%
%  Regel: Neuer Foliensatz  ->  neue Seite (\clearpage).
%  Pro Kapitel: Kernidee, Definitionen, Klausurtipps, Stolperfallen.
%
%  Kompilieren:  pdflatex INF4_1_Learnings.tex   (2x für Verweise)
% ============================================================
\documentclass[11pt,a4paper]{article}
\usepackage[utf8]{inputenc}
\usepackage[T1]{fontenc}
\usepackage[ngerman]{babel}
\usepackage{amsmath,amssymb}
\usepackage[table]{xcolor}
\usepackage[margin=18mm]{geometry}
\usepackage{booktabs}
\usepackage{array}
\usepackage{longtable}
\newcolumntype{L}[1]{>{\raggedright\arraybackslash}p{#1}}
\usepackage{enumitem}
\usepackage{tcolorbox}
\usepackage{fancyhdr}
\usepackage{lastpage}
\usepackage{fancyvrb}
\tcbuselibrary{skins,breakable}
\setlist{nosep,leftmargin=1.4em,topsep=2pt}

% ---- Farben -------------------------------------------------
\definecolor{ohmred}{RGB}{197,30,40}
\definecolor{defblue}{RGB}{30,90,160}
\definecolor{tipgreen}{RGB}{30,130,70}
\definecolor{pitred}{RGB}{180,40,40}

% ---- Kopf-/Fußzeile -----------------------------------------
\pagestyle{fancy}
\fancyhf{}
\lhead{\small\color{ohmred}\bfseries INF4/1 — Algorithmen \& Datenstrukturen}
\rhead{\small Learnings}
\cfoot{\small Seite \thepage\ von \pageref{LastPage}}
\renewcommand{\headrulewidth}{0.4pt}

% ---- Kapitel-Überschrift ------------------------------------
\newcommand{\kapitel}[2]{%
  \noindent{\color{ohmred}\rule{\linewidth}{2pt}}\par
  \vspace{2pt}
  {\LARGE\bfseries\color{ohmred} #1}\par
  {\large\color{black!70} #2}\par
  \vspace{4pt}
  {\color{ohmred}\rule{\linewidth}{0.5pt}}\par
  \vspace{8pt}
}

% ---- Boxen --------------------------------------------------
\newtcolorbox{defbox}[1]{colback=defblue!6,colframe=defblue,
  fonttitle=\bfseries,title={#1},boxrule=0.6pt,arc=2pt,left=4pt,right=4pt,top=3pt,bottom=3pt}
\newtcolorbox{tipbox}[1]{colback=tipgreen!7,colframe=tipgreen,
  fonttitle=\bfseries,title={#1},boxrule=0.6pt,arc=2pt,left=4pt,right=4pt,top=3pt,bottom=3pt}
\newtcolorbox{pitbox}[1]{colback=pitred!6,colframe=pitred,
  fonttitle=\bfseries,title={#1},boxrule=0.6pt,arc=2pt,left=4pt,right=4pt,top=3pt,bottom=3pt}

% ============================================================
\begin{document}
\thispagestyle{fancy}

% ---- Übersichtsseite: alle Algorithmen auf einen Blick ------
\kapitel{Algorithmen auf einen Blick}{Wie jedes Verfahren funktioniert — anschaulich erklärt, mit durchgerechnetem Mini-Beispiel}

\small
\renewcommand{\arraystretch}{1.25}
\begin{longtable}[l]{@{}L{2.5cm}L{7.4cm}L{6.6cm}@{}}
\toprule
\textbf{Algorithmus} & \textbf{So funktioniert es} & \textbf{Beispiel Schritt für Schritt} \\
\midrule
\endhead
\multicolumn{3}{@{}l}{\color{ohmred}\bfseries Grundlagen} \\
\addlinespace[2pt]
Euklid-ggT &
Sucht die größte Zahl, die zwei Zahlen ohne Rest teilt. Trick: der ggT ändert sich \emph{nicht}, wenn man die kleinere von der größeren Zahl abzieht. Das wiederholt man, bis beide Zahlen gleich sind — dieser Wert ist der ggT. &
$\mathrm{ggT}(12,8)$:\newline $12{-}8{=}4\Rightarrow(4,8)$\newline $8{-}4{=}4\Rightarrow(4,4)$\newline beide gleich $\Rightarrow\mathbf{4}$ \\
\addlinespace[3pt]
Max-Abschnitts\-summe \newline\textnormal{\footnotesize(Sliding Window)} &
Sucht das zusammenhängende Teilstück mit der größten Summe. Einmal von links durchgehen und aufsummieren; wird die laufende Summe negativ, hat der bisherige Anlauf nur geschadet $\to$ bei $0$ neu beginnen. Die größte je gesehene Summe ist die Antwort. &
$[\,3,\,-4,\,5,\,-1,\,2\,]$\newline $3$; dann $3{-}4{=}-1<0\Rightarrow$ Neustart\newline $5;\;5{-}1{=}4;\;4{+}2=\mathbf{6}$\newline bestes Stück $5,-1,2$ \\
\addlinespace[3pt]
\multicolumn{3}{@{}l}{\color{ohmred}\bfseries Sortieren} \\
\addlinespace[2pt]
Selection Sort &
Sucht in jedem Durchlauf das \emph{kleinste} Element des noch unsortierten Rests und tauscht es an dessen Anfang. Der sortierte Bereich wächst so von links um je ein Element. &
$[4,1,3,2]$\newline kleinstes $=1$ nach vorn: $[\mathbf{1},4,3,2]$\newline kleinstes des Rests $=2$: $[1,\mathbf{2},3,4]$\newline Rest schon sortiert $\Rightarrow$ fertig \\
\addlinespace[3pt]
Insertion Sort &
Wie Karten auf der Hand sortieren: nimm das nächste Element und schiebe es an die passende Stelle im bereits sortierten linken Teil; größere rücken dabei nach rechts. &
\underline{Teil} $=$ sortiert:\newline $[\underline{5}\mid 2,4,1]\to[\underline{2,5}\mid 4,1]$\newline $\to[\underline{2,4,5}\mid 1]\to[\underline{1,2,4,5}]$ \\
\addlinespace[3pt]
Bubble Sort &
Geht in Durchläufen über alle benachbarten Paare und tauscht, wenn links größer als rechts ist. Pro Durchlauf wandert das jeweils Größte ganz nach hinten (es „blubbert`` hoch). &
$[5,1,4,2]$\newline Durchlauf 1: $[1,4,2,\mathbf{5}]$ (5 hinten)\newline Durchlauf 2: $[1,2,4,5]$ \\
\addlinespace[3pt]
Merge Sort &
Teilt das Feld immer wieder in zwei Hälften, bis nur Einzelstücke bleiben. Dann fügt es paarweise zusammen: von zwei sortierten Teilen jeweils das kleinere vordere Element zuerst nehmen. &
$[5,2,4,1]$\newline teilen: $[5][2]\;[4][1]$\newline mischen: $[2,5]\;[1,4]$\newline mischen: $[1,2,4,5]$ \\
\addlinespace[3pt]
Quicksort &
Wählt ein Pivot-Element und sortiert alles Kleinere nach links, alles Größere nach rechts — das Pivot steht danach endgültig richtig. Beide Seiten werden nach demselben Prinzip weiter sortiert. &
$[4,2,7,1]$, Pivot $=4$\newline teilen: $[2,1]\;\mathbf{4}\;[7]$\newline links: $[2,1]\to[1,2]$\newline $\Rightarrow[1,2,4,7]$ \\
\addlinespace[3pt]
Heapsort &
Baut aus dem Feld einen Max-Heap (jeder Elternknoten $\ge$ seine Kinder), sodass das Größte oben steht. Nimm es nach hinten, stelle die Heap-Eigenschaft wieder her, wiederhole — das Feld fällt von hinten sortiert heraus. &
$[3,1,7,2]$\newline Max-Heap: $\mathbf{7}$ nach oben\newline $7$ ans Ende, Rest reparieren\newline $\dots\Rightarrow[1,2,3,7]$ \\
\addlinespace[3pt]
Counting Sort &
Vergleicht \emph{nicht}, sondern zählt: wie oft kommt jeder Wert vor? Danach gibt es die Werte der Größe nach so oft aus, wie sie gezählt wurden. Nur bei kleinem, bekanntem Wertebereich sinnvoll. &
$[2,0,2,1]$\newline zählen: $0{:}1,\;1{:}1,\;2{:}2$\newline ausgeben: $0,\,1,\,2,\,2$ \\
\pagebreak
\multicolumn{3}{@{}l}{\color{ohmred}\bfseries Bäume} \\
\addlinespace[2pt]
Binärer Suchbaum (BST) &
Jeder Knoten: links hängt alles $\le$, rechts alles $\ge$. Suchen/Einfügen: ab der Wurzel vergleichen und links/rechts absteigen; bei NIL gefunden bzw.\ als Blatt einfügen. Inorder-Walk (links–Knoten–rechts) gibt alles sortiert aus. &
einfügen $5,3,7,6$:\newline $5$ Wurzel; $3{<}5\to$ links;\newline $7{\ge}5\to$ rechts;\newline $6$: rechts von $5$, links von $7$ \\
\addlinespace[3pt]
BST: Löschen (2 Kinder) &
Nachfolger $=$ kleinster Wert im \emph{rechten} Teilbaum (dort immer links absteigen). Sein rechtes Kind rückt an seine alte Stelle, er selbst übernimmt Platz \emph{und beide Kinder} des Gelöschten — „der Nachfolger zieht allein um``. &
lösche Wurzel $8$, rechts $10{-}14$:\newline ab $10$ links absteigen $\to 10$\newline $10$ wird Wurzel, $14$ rückt nach\newline Rest bleibt unberührt \\
\addlinespace[3pt]
AVL-Baum: Einfügen &
BST, der sich selbst balanciert: nach jedem Einfügen/Löschen vom Elternknoten zur Wurzel die
Balancefaktoren prüfen ($BF=$ Höhe rechts $-$ Höhe links). Wird einer $\pm 2$, am \emph{untersten}
verletzten Knoten rotieren — Knoten und Kind in gleicher Richtung schief $\to$ einfache Rotation,
Knick $\to$ Doppelrotation. &
Kette $20{-}10{-}5$:\newline $BF(20){=}{-}2$, $BF(10){=}{-}1\le 0$\newline $\to$ Rechtsrotation an $20$:\newline $10$ hoch, $20$ wird \emph{rechtes} Kind \\
\addlinespace[3pt]
AVL: Doppel\-rotation &
Nötig, wenn der zu tiefe Teilbaum \emph{innen} liegt (Vorzeichen von Knoten und Kind entgegengesetzt
$=$ „Knick``). Erst das \emph{Kind} in die Gegenrichtung rotieren — das begradigt den Knick —, dann
den verletzten Knoten wie bei der einfachen Rotation. &
einfügen $50,30,70,60,65$:\newline $BF(70){=}{-}2$, $BF(60){=}{+}1$ (Knick)\newline $60$ nach links, dann $70$ nach rechts\newline $\Rightarrow 65$ oben, Kinder $60$ und $70$ \\
\addlinespace[3pt]
B-Baum: Einfügen &
Suchbaum mit dicken Knoten (Ordnung $m$: max.\ $2m{-}1$ Schlüssel). Neuer Schlüssel kommt \emph{immer} ins Blatt; \emph{volle} Knoten werden schon beim Abstieg gesplittet — der mittlere Schlüssel wandert in den Elternknoten. Nur ein Wurzel-Split macht den Baum höher. &
$m{=}2$, Blatt $[3\,5\,8]$ voll, füge $6$:\newline Split: $5$ hoch zu den Eltern,\newline Blätter $[3]$ und $[8]$\newline $6$ ins passende Blatt: $[6\,8]$ \\
\addlinespace[3pt]
B-Baum: Löschen (Unterlauf) &
Rutscht ein Knoten unter $m{-}1$ Schlüssel (Unterlauf), wird \emph{repariert}, nie gelöscht: kann ein Geschwister abgeben $\to$ \textbf{Verschieben} über den Elternknoten (Eltern-Schlüssel runter, Geschwister-Schlüssel rauf); sonst \textbf{Verschmelzen} mit Geschwister $+$ Trennschlüssel der Eltern. &
$m{=}2$: Blatt leer, Geschwister $[3\,5]$\newline kann abgeben $\to$ Verschieben:\newline Eltern-$10$ runter ins Blatt,\newline $5$ hoch zu den Eltern \\
\bottomrule
\end{longtable}

\vspace{6pt}
{\footnotesize Fett markiert steht ein Element endgültig · \underline{unterstrichen} $=$ schon sortierter Teil.
Komplexität, Stabilität und in-place stehen bei jedem Verfahren auf der jeweiligen Kapitelseite.\par
\vspace{2pt}
\textbf{Farblegende der Folgeseiten:} \textcolor{defblue}{\textbf{Blau}} $=$ Definitionen/Fakten ·
\textcolor{tipgreen}{\textbf{Grün}} $=$ Klausurtipps · \textcolor{pitred}{\textbf{Rot}} $=$ Stolperfallen.}

\clearpage
% ============================================================
%  KAP 01 — GRUNDBEGRIFFE
% ============================================================
\kapitel{01 — Allgemeines}{Grundbegriffe: Algorithmus, Datenstruktur, Darstellung}

\begin{defbox}{Algorithmus — Definition \& die 4 Merkmale}
Ein Algorithmus ist eine \textbf{endliche Folge} von Regeln, nach denen sich nach
\textbf{endlich vielen}, \textbf{eindeutig festgelegten} Schritten die Lösung einer
\textbf{Klasse von Problemen} ergibt.
\vspace{3pt}
\renewcommand{\arraystretch}{1.3}
\begin{tabular}{@{}p{3.0cm}p{4.3cm}p{6.5cm}@{}}
\textbf{Merkmal} & \textbf{Bedeutung} & \textbf{verletzt, wenn …} \\
\hline
Finitheit & endlich beschreibbar (endl. viele Regeln) & Anleitung unendlich lang \\
Terminiertheit & endet nach endlicher Zeit & Endlosschleife (Sprung zurück) \\
Determiniertheit & kein Interpretationsspielraum & vage Begriffe („angemessen``, „schön``) \\
Allgemeingültigkeit & Lösung für versch. Eingaben & nur ein fest verdrahteter Fall \\
\end{tabular}
\end{defbox}

\begin{itemize}
  \item \textbf{Datenstruktur:} Art und Weise, wie die von einem Algorithmus bearbeiteten
        Daten organisiert/vorgehalten werden. Algorithmus \& DS \textbf{bedingen sich
        gegenseitig}; manchmal ist das Erzeugen der DS aufwändiger als der Algorithmus
        selbst (Spaghetti-Sort).
  \item \textbf{Darstellungsformen} (informell $\to$ maschinennah): natürliche Sprache $\to$
        UML-Aktivitätsdiagramm $\to$ \textbf{Struktogramm} (Nassi-Shneiderman, DIN 66261)
        $\to$ \textbf{Pseudocode} (semi-formal) $\to$ \textbf{Implementierung} (Programmcode).
  \item \textbf{Zwei zentrale Eigenschaften:} \textbf{Korrektheit} (löst die Aufgabe richtig)
        und \textbf{Effizienz} (wenige Schritte + wenig Speicher).
  \item \textbf{Korrektheitsnachweis} (am Euklid-ggT): \emph{Terminierung} (Werte werden
        kleiner, Abbruch bei $x\le 0$) + \emph{Schleifeninvariante}
        $\mathrm{ggT}(x_t,y_t)=\mathrm{ggT}(x_{t+1},y_{t+1})$ $\to$ am Ende $x=y=$ ggT.
  \item \textbf{Abstraktes Maschinenmodell RAM} (Random Access Machine): $\infty$ viele
        Speicherzellen (je ein Integer bel. Größe); Befehle = lesen/schreiben, vergleichen,
        springen. \emph{Laufzeitkomplexität} = Anzahl Befehle, \emph{Speicherkomplexität}
        = Anzahl Zellen.
\end{itemize}

\begin{tipbox}{Klausurtipps}
\begin{itemize}
  \item Klassiker: \textbf{die 4 Merkmale nennen} oder „welche Eigenschaft ist hier verletzt?``
        $\to$ Eselsbrücke \textbf{F}initi–\textbf{T}ermini–\textbf{D}etermini–\textbf{A}llgemein.
  \item Korrektheit \& Effizienz sind \textbf{unabhängig von Programmiersprache/Implementierung}.
  \item RAM-Modell als Begründung, warum man \emph{zählt} statt \emph{misst}.
\end{itemize}
\end{tipbox}

\begin{pitbox}{Stolperfallen}
\begin{itemize}
  \item Finitheit (endliche \emph{Beschreibung}) $\neq$ Terminiertheit (endliche \emph{Laufzeit}) — nicht verwechseln!
  \item Pseudocode/Struktogramm ist \emph{kein} Programmcode: keine Implementierungsdetails (z.B. Tupel-Swap) verlangt.
\end{itemize}
\end{pitbox}

\clearpage
% ============================================================
%  KAP 02 — EFFIZIENZ / LANDAU-SYMBOLE
% ============================================================
\kapitel{02 — Effizienz von Algorithmen}{Komplexität \& Landau-Symbole}

\textbf{Kernidee.} Effizienz misst man als Funktion der Problemgröße $n$:
\begin{itemize}
  \item \textbf{Laufzeitkomplexität} = Anzahl ausgeführter Befehle
  \item \textbf{Speicherkomplexität} = Anzahl benutzter Speicherzellen
\end{itemize}
Für große $n$ zählt nur der \textbf{führende Term}; konstante Koeffizienten fallen weg:
\[ 1{,}5\,n^2 + 0{,}5\,n - 1 \;\longrightarrow\; \Theta(n^2)\ \text{(quadratisch)}. \]

\medskip
\textbf{Leitbeispiel „Maximale Abschnittssumme``} — eine Aufgabe, vier Komplexitätsklassen:

\begin{center}\small
\begin{tabular}{@{}llll@{}}
\toprule
\textbf{Ansatz} & \textbf{Idee} & \textbf{Komplexität} & \textbf{Name} \\
\midrule
MAXFOLGE1 & alle Teilfolgen, Summe neu berechnen & $O(n^3)$ & kubisch \\
MAXFOLGE2 & von jedem Start aufsummieren & $O(n^2)$ & quadratisch \\
MAXFOLGE3 & Teile-und-Herrsche (Rekursion) & $O(n\log n)$ & — \\
MAXFOLGE4 & einmal durchlaufen, bei Summe$<$0 neu & $O(n)$ & linear \\
\bottomrule
\end{tabular}
\end{center}
\textit{Praktischer Effekt bei $n=10^4$: kubisch $\approx$ 4738 s vs. linear $\approx$ 302 µs.}

\medskip
\begin{defbox}{Landau-Symbole (für alle $n \ge n_0$, positive Konstanten $c$)}
\renewcommand{\arraystretch}{1.4}
\begin{tabular}{@{}lll@{}}
$f \in O(g)$ & höchstens so schnell (\textbf{obere} Schranke) & $0 \le f(n) \le c\,g(n)$ \\
$f \in \Omega(g)$ & mindestens so schnell (\textbf{untere} Schranke) & $0 \le c\,g(n) \le f(n)$ \\
$f \in \Theta(g)$ & genauso schnell (\textbf{scharfe} Schranke) & $0 \le c_1 g(n) \le f(n) \le c_2 g(n)$ \\
\end{tabular}
\vspace{2pt}
\par\textbf{Beziehungen:}\quad
$\Theta(f) = O(f) \cap \Omega(f)$,\quad $\Theta(f) \subseteq O(f)$,\quad $\Theta(f) \subseteq \Omega(f)$.
\vspace{3pt}
\par\textbf{Mengenschreibweise} (denke $O,\Omega,\Theta$ als \emph{Mengen von Funktionen}):\par\nobreak\smallskip
\renewcommand{\arraystretch}{1.25}
\begin{tabular}{@{}l L{13.8cm}@{}}
$A \subseteq B$ & „$A$ ist \textbf{Teilmenge von oder gleich} $B$`` — jedes Element von $A$ liegt auch in $B$ (Gleichheit erlaubt, wie $\le$). \\
$A \not\subseteq B$ & „$A$ ist \textbf{keine} Teilmenge von $B$`` — es gibt ein Element in $A$, das \emph{nicht} in $B$ liegt. \\
$A \cap B$ & \textbf{Schnitt}: alles, was in $A$ \emph{und} $B$ liegt. \\
$A \cup B$ & \textbf{Vereinigung}: alles, was in $A$ \emph{oder} $B$ liegt. \\
\end{tabular}
\par\small Bsp.: $\Theta(f)\subseteq O(f)$ = „alles, was genau $f$-schnell ist, ist auch \emph{höchstens} $f$-schnell``.
\end{defbox}

\begin{tipbox}{Klausurtipps (laut Probeklausur)}
\begin{itemize}
  \item \textbf{Algorithmus entwerfen:} z.\,B. Summe von 3 aufeinanderfolgenden Zahlen in $O(n)$
        $\to$ \emph{Sliding Window}: erste 3 aufsummieren, dann pro Schritt vorne abziehen / hinten addieren.
  \item \textbf{Multiple Choice} zu $O/\Omega/\Theta$-Mengenbeziehungen — \emph{mit Minuspunkten} für Falsches!
        Richtig sind u.\,a. $\Theta(f)\subseteq O(f)$, $\Omega(f)\not\subseteq O(f)$, $\Theta(f)=O(f)\cap\Omega(f)$.
        Falsch: $\Theta(f)=O(f)\cup\Omega(f)$.
  \item \textbf{Komplexität ablesen:} Anzahl \emph{verschachtelter} Schleifen über $n$ = Exponent.
        2 verschachtelt $\to O(n^2)$, 3 $\to O(n^3)$, halbieren/rekursiv $\to \log n$.
\end{itemize}
\end{tipbox}

\begin{pitbox}{Stolperfallen}
\begin{itemize}
  \item $O$ ist \emph{nur} obere Schranke: auch $n \in O(n^2)$ ist korrekt! Für die \emph{exakte} Klasse $\Theta$ nehmen.
  \item Bei der MC-Aufgabe im Zweifel \emph{nicht} ankreuzen — Falsches kostet einen Punkt.
  \item Nicht verschachtelte Schleifen addieren sich nur ($O(n)+O(n)=O(n)$); nur \emph{verschachtelte} multiplizieren.
\end{itemize}
\end{pitbox}

\clearpage
% ============================================================
%  KAP 03 — ELEMENTARES SORTIEREN
% ============================================================
\kapitel{03 — Elementares Sortieren}{Selection, Insertion, Bubble Sort + Eigenschaften}

\textbf{Kernidee.} Ein Sortierverfahren bringt ein Feld in eine Ordnung. Bewertet wird es
nach drei Eigenschaften, die in der Klausur abgefragt werden:

\begin{defbox}{Drei Eigenschaften eines Sortierverfahrens}
\renewcommand{\arraystretch}{1.3}
\begin{tabular}{@{}p{2.6cm}p{11.6cm}@{}}
\textbf{korrekt} & terminiert (finit) \emph{und} liefert am Ende eine sortierte Permutation der Eingabe \\
\textbf{in-place} & Zusatzspeicher \textbf{wächst nicht mit $n$} $\Rightarrow O(1)$ (ein paar Hilfsvariablen erlaubt) \\
\textbf{stabil} & gleichwertige Elemente behalten ihre \emph{relative} Reihenfolge über die ganze Laufzeit \\
\end{tabular}
\end{defbox}

\begin{defbox}{Die drei elementaren Verfahren (alle in-place)}
\renewcommand{\arraystretch}{1.3}
\small
\begin{tabular}{@{}p{2.3cm}cccp{6.3cm}@{}}
\textbf{Verfahren} & \textbf{Best} & \textbf{Avg} & \textbf{Worst} & \textbf{Mechanik} \\
\hline
Selection & $n^2$ & $n^2$ & $n^2$ & sucht Minimum des unsortierten Rests, hängt es an den sortierten Teil an \\
Insertion & $\boldsymbol{n}$ & $n^2$ & $n^2$ & nimmt 1. unsort. Element, schiebt es von hinten an die richtige Stelle des sort. Teils \\
Bubble & $\boldsymbol{n}$ & $n^2$ & $n^2$ & vertauscht benachbarte Paare, „blubbert`` das Größte nach oben; mehrere Durchläufe \\
\end{tabular}
\vspace{2pt}
\par\textbf{stabil:} Insertion \checkmark, Bubble \checkmark, Selection \textbf{nein}.\quad
\textbf{Swaps:} Selection nur $O(n)$, Bubble bis $O(n^2)$.
\end{defbox}

\textbf{Best-/Worst-/Average-Case allgemein:}
Best = Feld bereits sortiert · Worst = umgekehrt sortiert (allg.: jedes Element muss mit
jedem verglichen werden) · Average = zufällig.

\begin{defbox}{Korrektheitsnachweis über Schleifeninvariante}
\textbf{Invariante} (= „un-veränderlich``): eine Aussage, die \emph{vor jedem} Schleifendurchlauf
wahr ist — egal in welcher Runde. Wie ein \emph{Versprechen, das die Schleife nie bricht},
während sich Index und Feldinhalt drumherum ändern. Korrektheit zeigt man in \textbf{drei Schritten}:
\begin{enumerate}
  \item \textbf{Initialisierung:} Die Invariante gilt vor dem 1. Durchlauf.
  \item \textbf{Aufrechterhaltung:} Gilt sie vor einem Durchlauf, so gilt sie auch danach (vor dem nächsten).
  \item \textbf{Terminierung:} Die Schleife endet, und aus der Invariante folgt die gewünschte Korrektheitsaussage.
\end{enumerate}
\textit{Beispiel Insertion Sort} — Versprechen: „Der linke Teil $A[1..\,j{-}1]$ ist immer schon
\emph{sortiert} (und enthält die ursprünglichen Elemente).``
\begin{itemize}
  \item \emph{Init} ($j{=}2$): $A[1..1]$ ist ein Element $\to$ trivial sortiert. \checkmark
  \item \emph{Erhaltung:} Der Durchlauf schiebt $A[j]$ an die richtige Stelle $\to$ $A[1..\,j]$ sortiert.
  \item \emph{Terminierung} ($j{=}n{+}1$): der „linke Teil`` ist das ganze Feld $A[1..\,n]$ $\to$ sortiert. \checkmark
\end{itemize}
Durchgespielt an $[5,2,4,1]$: $[\underline{5}]\,2\,4\,1 \to [\underline{2,5}]\,4\,1 \to [\underline{2,4,5}]\,1 \to [\underline{1,2,4,5}]$
— der unterstrichene linke Teil bleibt in jedem Schritt sortiert.
(Gleiches Muster wie Euklid-ggT in Kap. 01: Terminierung + Invariante.)
\end{defbox}

\begin{defbox}{Korrektheit der Rekursion: (strukturelle) Induktion — aus Praktikum 02}
Bei \emph{rekursiven} Verfahren reicht die Schleifeninvariante nicht: sie beweist nur einen
Teilschritt. Die Rekursion drumherum zeigt man per \textbf{Induktion über die Eingabegröße $n$}
(Dominoprinzip), wobei Rekursion und Induktion dieselbe Form haben:
\begin{itemize}
  \item \textbf{Basisfall:} kleinste Eingabe ($n{=}1$) ist trivial korrekt = Abbruchbedingung der Rekursion.
  \item \textbf{Annahme:} die rekursiven Aufrufe für \emph{kleinere} Teile ($<n$) arbeiten bereits korrekt.
  \item \textbf{Schritt:} aus den korrekten Teil-Ergebnissen folgt das korrekte Gesamt-Ergebnis.
\end{itemize}
\textit{Bsp. Merge Sort (Kap. 04):} Annahme liefert zwei sortierte Hälften $\to$ der (per
Invariante bewiesene) Merge fügt sie zu einem sortierten Feld $\to$ Länge $n$ sortiert.
\par\textbf{Kernaussage:} \emph{Schleifeninvariante} (innerer Schritt) + \emph{Induktion}
(Rekursion drumherum) ergeben \textbf{zusammen} den vollständigen Korrektheitsnachweis.
\end{defbox}

\begin{tipbox}{Klausurtipps (laut Probeklausur)}
\begin{itemize}
  \item \textbf{SELECT-Sort von Hand tracen} (5 P): Feld \emph{nach jedem Tausch} aufschreiben.
        Vorgehen: Min im Rest suchen, mit erstem unsortierten Element tauschen, Grenze um 1 nach rechts.
  \item Verfahren \textbf{einordnen}: stabil ja/nein, in-place ja/nein, Best/Avg/Worst nennen.
  \item Best-Case \textbf{begründen}, nicht nur nennen (Insertion/Bubble brechen bei sortiert früh ab).
\end{itemize}
\end{tipbox}

\begin{pitbox}{Stolperfallen}
\begin{itemize}
  \item \textbf{in-place $=O(1)$ unabhängig von $n$} — nicht „Variable ja, Array nein``! Ein Hilfsarray
        \emph{fester} Größe wäre auch in-place; $n$ einzelne Variablen wären es \emph{nicht}.
        \textbf{Merge Sort} braucht ein Hilfsarray der Größe $n$ $\Rightarrow O(n)$ $\Rightarrow$ \emph{nicht} in-place.
  \item \textbf{Selection Sort hat keinen guten Best-Case:} auch bei sortiertem Feld $n^2$, weil das
        Min-Suchen \emph{immer} den ganzen Rest durchläuft (datenunabhängig, $\approx n^2/2$ Vergleiche).
        Die Regel „Best = sortiert $\Rightarrow n$`` gilt \emph{nur} für Insertion/Bubble!
  \item \textbf{Selection Sort ist meist nicht stabil:} der Min-Swap kann Gleiche überspringen,
        z.\,B. $[5_a,5_b,1]\to[1,5_b,5_a]$.
  \item \textbf{Bubble-Best $n$ nur mit „swapped``-Flag:} ohne Abbruch-Optimierung (kein Tausch im
        Durchlauf $\Rightarrow$ fertig) läuft naives Bubble immer $n^2$.
\end{itemize}
\end{pitbox}

\clearpage
% ============================================================
%  KAP 04 — FORTGESCHRITTENES SORTIEREN
% ============================================================
\kapitel{04 — Fortgeschrittenes Sortieren}{Quicksort, Heapsort, untere Schranke, Counting Sort}

\textbf{Kernidee.} Über die elementaren $\Theta(n^2)$-Verfahren hinaus: zwei
\emph{vergleichsbasierte} $n\log n$-Verfahren (Quicksort, Heapsort) und — wenn der Wertebereich
bekannt ist — ein \emph{nicht} vergleichsbasiertes Linearverfahren (Counting Sort).

\begin{defbox}{Überblick ($k$ = Größe des Wertebereichs bei Counting Sort)}
\renewcommand{\arraystretch}{1.3}\small
\begin{tabular}{@{}p{2.4cm}ccccp{5.4cm}@{}}
\textbf{Verfahren} & \textbf{Best} & \textbf{Avg} & \textbf{Worst} & \textbf{stabil / in-place} & \textbf{Idee} \\
\hline
Quicksort & $n\log n$ & $n\log n$ & $\boldsymbol{n^2}$ & nein / ja ($O(\log n)$ Stack) & Pivot $\to$ partitionieren $\to$ rekursiv \\
Heapsort & $n\log n$ & $n\log n$ & $n\log n$ & nein / ja & Max-Heap bauen, Wurzel abtragen \\
Counting Sort & $n{+}k$ & $n{+}k$ & $n{+}k$ & \textbf{ja} / nein & zählen statt vergleichen \\
\end{tabular}
\end{defbox}

\begin{defbox}{Quicksort — Rezept}
\begin{enumerate}
  \item \textbf{Pivot wählen} (oft \textbf{Median-of-Three} aus erstem/mittlerem/letztem Element) und
        an die \textbf{letzte Stelle} tauschen.
  \item \textbf{Zwei-Zeiger-Partition:} linker Zeiger $L$ läuft nach rechts bis Wert $\ge$ Pivot,
        rechter Zeiger $R$ nach links bis Wert $\le$ Pivot — \textbf{beide vergleichen mit dem Pivot,
        nicht miteinander}. Stehen beide $\Rightarrow$ \textbf{Tausch}.
  \item \textbf{Kreuzen} sich die Zeiger ($L$ rechts von $R$) $\Rightarrow$ Pivot an die Position von
        $L$ tauschen — jetzt \textbf{endgültig platziert}.
  \item \textbf{Rekursion} auf linkem und rechtem Teil, jeweils \textbf{ohne} den Pivot.
\end{enumerate}
\textbf{Worst Case $n^2$} bei ungünstigem Pivot (immer kleinstes/größtes Element, z.\,B. bereits
sortiertes Feld). Median-of-Three entschärft genau das.
\end{defbox}

\begin{defbox}{Quicksort-Trace (jeweils \emph{Moment vor dem Tausch})}
\small\textbf{Pass 1} — Feld \texttt{4 2 7 1 6 3 8 5}, \textbf{Pivot = 6} (an letzte Stelle $\to$ 5 rutscht auf Idx 4):
\begin{Verbatim}[fontsize=\small,xleftmargin=6pt]
  4  2  7  1  5  3  8  6     L bei 7 (>=6), R bei 3 (<=6)
        L        R     P     -> L < R:  Tausch 7 <-> 3

  4  2  3  1  5  7  8  6     weiter: L bei 7 (Idx5), R bei 5 (Idx4)
              R  L     P     -> gekreuzt:  Pivot setzen  7 <-> 6

  Ergebnis:  [ 4 2 3 1 5 ]  6  [ 8 7 ]     (6 endgueltig platziert)
\end{Verbatim}
\textbf{Pass 2} — linkes Teilfeld \texttt{4 2 5 1 3}, \textbf{Pivot = 3} (schon letzte Stelle):
\begin{Verbatim}[fontsize=\small,xleftmargin=6pt]
  4  2  5  1  3     L bei 4 (>=3), R bei 1 (<=3)
  L        R  P     -> L < R:  Tausch 4 <-> 1

  1  2  5  4  3     weiter: L bei 5 (Idx2), R bei 2 (Idx1)
     R  L     P     -> gekreuzt:  Pivot setzen  5 <-> 3

  Ergebnis:  [ 1 2 ]  3  [ 4 5 ]     -> nach Rekursion fertig sortiert
\end{Verbatim}
\end{defbox}

\begin{defbox}{Heapsort}
\begin{enumerate}
  \item Feld als \textbf{Max-Heap} aufbauen: Kind $\le$ Elternknoten. Als Array liegen die Kinder von
        Index $i$ bei $2i{+}1$ und $2i{+}2$.
  \item \textbf{Wurzel} (= Maximum) mit dem \emph{letzten} Heap-Element tauschen und die Heap-Größe um
        1 \textbf{verkleinern} — das Max liegt jetzt endgültig hinten (``löschen'' = nur Größe$-1$).
  \item \texttt{heapify}/\texttt{sift-down} von der Wurzel $\Rightarrow$ Max-Heap wiederhergestellt.
  \item Schritte 2–3 wiederholen, bis der Heap leer ist.
\end{enumerate}
$\Theta(n\log n)$ in \emph{allen} Fällen, \textbf{in-place}, aber \textbf{nicht stabil}
(Heap-Aufbau ist $O(n)$).
\end{defbox}

\begin{defbox}{Untere Schranke \& Counting Sort}
\textbf{Untere Schranke.} Jedes \textbf{vergleichsbasierte} Sortierverfahren braucht im Worst Case
$\Omega(n\log n)$ Vergleiche: der Entscheidungsbaum hat $n!$ Blätter (alle Permutationen) $\Rightarrow$
Höhe $\ge \log_2(n!) = \Theta(n\log n)$. \emph{Schneller als $n\log n$ geht mit reinen Vergleichen nicht.}
\vspace{4pt}
\par\textbf{Counting Sort} umgeht die Schranke, weil es \textbf{nicht vergleicht, sondern zählt}
(Wertebereich $0..k$ muss bekannt \& klein sein):
\begin{enumerate}
  \item \textbf{Zählen:} Hilfsarray $C$ — wie oft kommt jeder Wert vor.
  \item \textbf{Präfixsumme} über $C$: danach ist $C[v]$ = Anzahl Elemente $\le v$ (= Endposition von $v$).
  \item \textbf{Von hinten} durch die Eingabe (das sichert \textbf{Stabilität}): Element $v$ an
        Ausgabeposition $C[v]-1$ schreiben, dann $C[v]$ um 1 verringern.
\end{enumerate}
$O(n+k)$, \textbf{stabil}, \textbf{nicht in-place} (zweites Feld nötig).
\end{defbox}

\begin{tipbox}{Klausurtipps}
\begin{itemize}
  \item \textbf{Quicksort/Heapsort von Hand tracen} — Zwischenzustände angeben (wie beim SELECT-Sort, §Kap.03).
        Bei Quicksort: Pivot markieren, nach jedem Tausch das Feld notieren.
  \item Verfahren \textbf{einordnen}: Best/Avg/Worst, stabil, in-place — besonders Quicksorts
        \emph{Worst Case $n^2$} und Heapsorts \emph{immer $n\log n$}.
  \item Frage nach der \textbf{$n\log n$-Grenze}: gilt nur für \emph{vergleichsbasierte} Verfahren;
        Counting Sort ist die Ausnahme (setzt bekannten Wertebereich voraus).
\end{itemize}
\end{tipbox}

\begin{pitbox}{Stolperfallen}
\begin{itemize}
  \item \textbf{Quicksort-Zeiger vergleichen mit dem \emph{Pivot}}, nicht miteinander: $L$ stoppt bei
        $\ge$ Pivot, $R$ bei $\le$ Pivot; erst \emph{dann} wird getauscht.
  \item \textbf{$\Omega(n\log n)$ ist eine \emph{untere} Schranke:} vergleichsbasierte Sortierer kommen
        \emph{nicht darunter} (nicht besser als $n\log n$) — nicht „erreichen nur``.
  \item \textbf{Counting-Sort-Zählarray wird am Ende \emph{nicht} komplett null:} jedes $C[v]$ endet bei
        „Anzahl Elemente \emph{echt kleiner} als $v$``. Fertig ist man, wenn alle $n$ Eingabe-Elemente
        einsortiert sind.
  \item \textbf{Heapsort ist \emph{nicht} stabil}, obwohl in-place — gleiche Werte können ihre
        Reihenfolge verlieren.
\end{itemize}
\end{pitbox}

\clearpage
% ============================================================
%  KAP 06 — BINÄRE SUCHBÄUME
% ============================================================
\kapitel{06 — Binäre Suchbäume (BST)}{Binärbaum, Suchbaum-Eigenschaft, Suchen / Einfügen / Löschen}

\textbf{Kernidee.} Statt eine Menge erst unsortiert in ein Array zu legen und später zu
sortieren, wird sie schon \emph{beim Einfügen} in eine Baumstruktur einsortiert, die
\textbf{effizientes Suchen} erlaubt und \textbf{dynamisch wächst} (Anzahl vorab unbekannt).

\begin{defbox}{Binärbaum — Begriffe}
\begin{itemize}
  \item Jeder Knoten hat \textbf{höchstens zwei} Kinder (daher \emph{binär}); jeder Knoten
        außer der \textbf{Wurzel} hat genau einen Elternknoten.
  \item Knoten \textbf{ohne} Kinder $=$ \textbf{äußere Knoten} $=$ \textbf{Blätter}; die übrigen $=$ \emph{innere} Knoten.
  \item \textbf{Höhe $h$} $=$ längster Weg von der Wurzel zu einem Blatt, gezählt in
        \textbf{Knoten-Ebenen — die Wurzel zählt mit}: Kette aus $n$ Knoten $\Rightarrow h=n$
        (Konvention laut AVL-Folien, Kap. 07).
  \item Anders als beim Heap kann ein Binärbaum \textbf{sehr unbalanciert} sein $\Rightarrow$
        \emph{kein} einfaches Abbilden auf ein Array; realisiert über \textbf{verkettete} Knoten
        (Zeiger \texttt{left}, \texttt{right}).
\end{itemize}
\end{defbox}

\begin{defbox}{Binäre-Suchbaum-Eigenschaft (BST)}
Ein Binärbaum ist ein \textbf{Suchbaum}, wenn die Knoten vergleichbare Schlüssel (\emph{keys})
tragen und für \emph{jeden} Knotenwert $x$ gilt:
\[ \forall\,v \text{ im \textbf{linken} Teilbaum}:\ v \le x
   \qquad
   \forall\,w \text{ im \textbf{rechten} Teilbaum}:\ x \le w \]
Die Bedingung gilt für den \textbf{gesamten Teilbaum}, nicht nur die direkten Kinder. Gleiche
Werte dürfen laut Folie auf \emph{beide} Seiten. Ein BST zu einer Menge ist \textbf{nicht
eindeutig} (die Einfügereihenfolge bestimmt die Form).
\end{defbox}

\begin{defbox}{Die drei Grund-Operationen}
\renewcommand{\arraystretch}{1.35}
\begin{tabular}{@{}p{2.5cm}p{11.5cm}@{}}
\textbf{Sortieren} & \textbf{Inorder-Tree-Walk}: links $\to$ Knoten ausgeben $\to$ rechts. Gibt die
  Werte \textbf{aufsteigend sortiert} aus (``inorder'', weil der Knoten \emph{zwischen} den
  Teilbäumen verarbeitet wird). Jeder Knoten genau einmal $\Rightarrow O(n)$. \\
\textbf{Suchen} & \textsc{Tree-Search}: mit Wurzel vergleichen; $key<x.key\to$ links, sonst rechts;
  bei NIL nicht enthalten. Laufzeit $O(h)$. \\
\textbf{Einfügen} & Von der Wurzel \textbf{absteigen} (wie beim Suchen) bis zu einem \textbf{NIL};
  dort das neue Element als \textbf{Blatt} einhängen. Laufzeit $O(h)$. \\
\end{tabular}
\end{defbox}

\begin{defbox}{Löschen — die drei Fälle (nach Foliensatz)}
\begin{enumerate}
  \item \textbf{Blatt} (kein Kind): einfach entfernen (auf NIL setzen) — die restliche Struktur bleibt unberührt.
  \item \textbf{Genau ein Kind:} den Zeiger auf den zu löschenden Knoten auf dessen
        \textbf{einziges (nicht-leeres) Kind umlenken} — der Teilbaum rückt hoch.
  \item \textbf{Zwei Kinder:} durch den \textbf{symmetrischen Nachfolger} ersetzen $=$
        \textbf{kleinster Wert im rechten Teilbaum} (im rechten Teilbaum immer nach links absteigen).
        \begin{itemize}
          \item Einfachster Fall: der Nachfolger ist ein \textbf{Blatt}.
          \item Ist er ein \emph{innerer} Knoten, hat er \textbf{kein linkes Kind} (sonst wäre \emph{das}
                der Nachfolger) $\Rightarrow$ sein \textbf{rechtes} Kind nachrücken lassen, dann den
                Nachfolger an die Lücke setzen.
        \end{itemize}
\end{enumerate}
\end{defbox}

\begin{defbox}{Lösch-Trace (2 Kinder) — Nachfolger aus dem rechten Teilbaum}
\small Baum, aufgebaut durch Einfügen von $50,30,70,20,40,60,80,35$:
\begin{Verbatim}[fontsize=\small,xleftmargin=6pt]
          50                Loesche 30  (hat zwei Kinder: 20 und 40)
        /    \              Nachfolger = kleinster Wert im RECHTEN
      30      70            Teilbaum von 30: ab 40 nach links -> 35
     /  \    /  \           (35 ist hier ein Blatt)
    20  40  60  80
        /
       35
\end{Verbatim}
$\Rightarrow$ 35 ersetzt die 30, das alte Blatt 35 verschwindet — Ergebnis ist wieder ein gültiger BST:
\begin{Verbatim}[fontsize=\small,xleftmargin=6pt]
          50
        /    \
      35      70
     /  \    /  \
    20  40  60  80
\end{Verbatim}
\end{defbox}

\begin{defbox}{Laufzeit-Komplexität der Operationen}
\renewcommand{\arraystretch}{1.25}
\begin{tabular}{@{}ll@{}}
Traversieren & $O(n)$ \\
Suchen / Einfügen / Löschen & $O(h)$ \\
\end{tabular}
\par\vspace{3pt} $h=$ Höhe des Baums: \textbf{balanciert} $h\approx\log_2 n \Rightarrow O(\log n)$;
\textbf{entartet} (z.\,B. sortiert eingefügt $\to$ lineare Kette) $h=n \Rightarrow O(n)$.
\par\vspace{2pt}\textbf{Herleitung über Rekurrenz (Praktikum 4):} entartet
$T(n)=T(n{-}1)+c \Rightarrow O(n)$; balanciert $T(n)=T(n/2)+c \Rightarrow O(\log n)$
(pro Schritt halbiert sich der Restbaum).
\emph{Beispiel $n=8$, sortiert eingefügt:} \texttt{search(8)} braucht \textbf{8} Vergleiche,
im ausgewogenen Baum nur \textbf{4}. Ausgewogen wird der BST nur bei günstiger
Einfügereihenfolge (Median zuerst, dann rekursiv die Median-Hälften) — \emph{garantieren}
kann $h=O(\log n)$ erst der AVL-Baum (Kap. 07).
\end{defbox}

\begin{tipbox}{Klausurtipps}
\begin{itemize}
  \item \textbf{Bäume von Hand} (8 P): Element einfügen bzw.\ löschen \emph{zeichnen}. Beim Löschen
        zuerst den \textbf{Fall bestimmen} (0 / 1 / 2 Kinder).
  \item \textbf{Inorder-Walk} liefert die sortierte Reihenfolge — ideal zum \emph{Nachprüfen}, ob dein
        gezeichneter Baum noch ein gültiger BST ist.
  \item \textbf{Inorder vs. Level-Order (Praktikum 4):} Inorder ist bei \emph{jeder} BST-Form gleich
        (sortiert!) und verrät nichts über die Gestalt — ob der Baum \emph{balanciert oder entartet}
        ist, sieht man erst am Level-Order-Walk (ebenenweise Ausgabe).
  \item Beim 2-Kinder-Löschen die \textbf{Folien-Methode} nehmen: Nachfolger $=$ kleinster Wert im
        \emph{rechten} Teilbaum.
\end{itemize}
\end{tipbox}

\begin{pitbox}{Stolperfallen}
\begin{itemize}
  \item \textbf{BST-Eigenschaft gilt für den ganzen Teilbaum}, nicht nur die direkten Kinder — auch
        ein Enkel im linken Teilbaum muss $\le x$ sein.
  \item \textbf{2 Kinder $\to$ Nachfolger aus dem \emph{rechten} Teilbaum} (kleinster Wert dort),
        \emph{nicht} das äußerste rechte Blatt des linken Teilbaums. (Der Vorgänger von links ergäbe
        ebenfalls einen gültigen BST, ist aber nicht die Folien-Konvention.)
  \item \textbf{Beim 2-Kinder-Löschen zieht NUR der Nachfolger um:} (1) sein rechtes Kind rückt an
        seine \emph{alte} Stelle, (2) er selbst übernimmt Position \textbf{und beide Kinder} des
        Gelöschten — der restliche Baum bleibt unberührt! \emph{Nicht} den übrigen Teilbaum unter
        das nachrückende Kind hängen (ergäbe zwar oft einen gültigen BST, aber nicht das Ergebnis
        des Verfahrens). Merksatz: \emph{„Der Nachfolger zieht allein um.``}
  \item \textbf{Sortierte Eingabe erzeugt einen entarteten Baum} (lineare Kette) $\Rightarrow h=n$,
        alle Operationen $O(n)$. Genau das behebt der AVL-Baum (Kap. 07).
\end{itemize}
\end{pitbox}

\clearpage
% ============================================================
%  KAP 07 — AVL-BÄUME
% ============================================================
\kapitel{07 — AVL-Bäume}{Balancefaktor, Rotationen, Einfügen und Löschen mit Rebalancieren}

\textbf{Kernidee.} Alle BST-Operationen kosten $O(h)$ — aber die Höhe hängt von der
Einfügereihenfolge ab: eine \emph{sortierte} Folge entartet zur Kette mit $h=n$ (das Problem
ist vergleichbar mit dem Worst Case von Quicksort). Ein \textbf{AVL-Baum} (nach
\textbf{A}delson-\textbf{V}elski und \textbf{L}andis) ist ein binärer Suchbaum, der bei
\emph{jedem} Einfügen oder Löschen „balanciert`` wird — so bleibt $h=O(\log n)$
\textbf{garantiert}, unabhängig von der Eingabereihenfolge.

\begin{defbox}{Balancefaktor (Folie 4)}
Für einen Knoten $k$ ist der Balancefaktor definiert als
\[ BF(k) \;=\; \mathrm{HEIGHT}(k.right) \;-\; \mathrm{HEIGHT}(k.left)
   \qquad \text{(\textbf{rechts minus links}!)} \]
und für \emph{jeden} Knoten des Baums muss gelten: \; $BF(k)\in\{-1,\,0,\,1\}$.
\begin{itemize}
  \item \textbf{Höhen-Konvention:} in \textbf{Knoten-Ebenen} gezählt — $\mathrm{HEIGHT}(\mathrm{NIL})=0$,
        Blatt $=1$, die Wurzel zählt mit (Kette aus $n$ Knoten $\Rightarrow h=n$).
  \item \textbf{Lesart:} $BF<0$ $=$ \textbf{linkslastig} · $BF>0$ $=$ \textbf{rechtslastig} · $\pm 2$ $=$ verletzt.
  \item \textbf{Knoten-Datenstruktur} (Folie 5): \texttt{key}, \texttt{balance}, \texttt{left},
        \texttt{right}, \texttt{parent} — der Eltern-Verweis ermöglicht das Aufsteigen beim Rebalancieren.
\end{itemize}
\end{defbox}

\begin{defbox}{Rotation — die Mechanik (Rechtsrotation; Linksrotation spiegelbildlich)}
\begin{Verbatim}[fontsize=\small,xleftmargin=6pt]
        k                    l          l steigt auf, k wird RECHTES Kind von l.
       / \                  / \         Der mittlere Teilbaum b wechselt die
      l   c      ->        a   k        Seite: war rechtes Kind von l,
     / \                      / \       wird linkes Kind von k.
    a   b                    b   c      Inorder a,l,b,k,c bleibt unveraendert!
\end{Verbatim}
\end{defbox}

\begin{defbox}{Welche Rotation? — Entscheidungstabelle (Folien 6--10)}
\renewcommand{\arraystretch}{1.35}\small
\begin{tabular}{@{}llp{6.6cm}@{}}
\textbf{Verletzung} & \textbf{Kind} & \textbf{Rotation} \\
\hline
$BF(k)=-2$ (linkslastig) & $BF(k.left)\le 0$ (gleiche Ri.) & \textbf{einfach}: Rechtsrotation an $k$ \\
$BF(k)=-2$ & $BF(k.left)>0$ (\textbf{Knick}) & \textbf{doppelt}: erst $k.left$ nach links, dann $k$ nach rechts \\
$BF(k)=+2$ (rechtslastig) & $BF(k.right)\ge 0$ (gleiche Ri.) & \textbf{einfach}: Linksrotation an $k$ \\
$BF(k)=+2$ & $BF(k.right)<0$ (\textbf{Knick}) & \textbf{doppelt}: erst $k.right$ nach rechts, dann $k$ nach links \\
\end{tabular}
\par\vspace{3pt}
\textbf{Merkregel:} Knoten und Kind in \emph{gleicher} Richtung schief (oder Kind $0$) $\to$ \textbf{eine}
Rotation in die Gegenrichtung der Schieflage. Kind in \emph{Gegenrichtung} schief („Knick``: der zu
tiefe Teilbaum liegt \emph{innen}) $\to$ \textbf{Doppelrotation} — die erste begradigt den Knick,
die zweite balanciert.
\par\vspace{2pt}
{\footnotesize (\textsc{Avl-Balance}, Folie 10, entscheidet äquivalent über Höhen: bei $-2$
Rechtsrotation, falls $\mathrm{HEIGHT}(k.left.left)\ge\mathrm{HEIGHT}(k.left.right)$, sonst
Doppelrotation — und steigt danach rekursiv zum Elternknoten auf, bis zur Wurzel.)}
\end{defbox}

\begin{defbox}{Rezept: Einfügen (Folie 11)}
\begin{enumerate}
  \item Neuen Knoten \textbf{wie im BST als Blatt} einfügen (Kap.-06-Seite).
  \item Vom \textbf{Elternknoten des neuen Blatts} aufsteigend die Balancefaktoren neu berechnen.
  \item Am \textbf{untersten} Knoten mit $BF=\pm 2$: Rotationstyp per Tabelle bestimmen, rotieren.
  \item Weiter bis zur Wurzel prüfen (beim Einfügen heilt die eine Rotation i.\,d.\,R. den ganzen Pfad).
  \item \textbf{Selbstcheck:} Inorder noch sortiert? Alle $BF\in\{-1,0,1\}$?
\end{enumerate}
\end{defbox}

\begin{defbox}{Rezept: Löschen (Folie 12 + Praktikum 4)}
\begin{enumerate}
  \item Knoten \textbf{wie im BST löschen} — 3 Fälle (Kap.-06-Seite); bei 2 Kindern:
        \emph{„der Nachfolger zieht allein um``}.
  \item Rebalancieren beginnt beim \textbf{Elternknoten des (physisch) entfernten Knotens} — beim
        2-Kinder-Fall also beim alten Elternknoten des \emph{Nachfolgers}. \emph{Warum dort?} Erst ab
        diesem Knoten haben sich Höhen geändert; der gelöschte Knoten selbst ist weg
        (Praktikum 4: \texttt{balance(parent)}, nicht \texttt{balance(node)}).
  \item Aufsteigend bis zur Wurzel prüfen — die Verletzung kann \textbf{weit oberhalb} der Löschstelle
        auftreten, und anders als beim Einfügen können \textbf{mehrere Rotationen} nötig werden.
\end{enumerate}
\end{defbox}

\begin{defbox}{Zwei Mini-Traces}
\small\textbf{Einfache Rotation} — $5$ in den linkslastigen Ast $20{-}10$ eingefügt:
\begin{Verbatim}[fontsize=\small,xleftmargin=6pt]
      20                    10        BF(20) = -2,  BF(10) = -1 <= 0
     /                     /  \       -> gleiche Richtung -> EINFACHE
   10           ->        5    20     Rechtsrotation an 20: 10 steigt
   /                                  auf, 20 wird ihr RECHTES Kind
  5
\end{Verbatim}
\textbf{Doppelrotation — durch Löschen ausgelöst} (Praktikum 4, nach Löschen von $1,3,10$):
das \textbf{Blatt 20} wird gelöscht — die Verletzung erscheint \emph{zwei Ebenen höher} bei der Wurzel:
\begin{Verbatim}[fontsize=\small,xleftmargin=6pt]
      12                12                  7         BF(12) = -2, BF(5) = +1 > 0
     /  \              /  \               /   \       -> Knick -> DOPPELROTATION:
    5    15    ->     7    15    ->      5     12     erst 5 nach links
   / \               / \                / \   /  \    (6 wechselt zu 5),
  4   7             5   8              4   6 8    15  dann 12 nach rechts
     / \           / \                                (8 wechselt zu 12)
    6   8         4   6
\end{Verbatim}
End-BFs: \emph{alle} $0$. Elternknoten $15$ der Löschstelle blieb gültig — deshalb \textbf{bis zur Wurzel} prüfen!
\end{defbox}

\begin{defbox}{Laufzeit-Komplexität (Folien 11--12)}
\renewcommand{\arraystretch}{1.25}
\begin{tabular}{@{}ll@{}}
Suchen & $O(\log n)$ — die Höhe ist garantiert logarithmisch \\
Balancieren (ein Aufstieg zur Wurzel) & $O(\log n)$ \\
Einfügen inkl.\ Balancieren & $O(\log n)$ \textbf{auch im Worst Case} \\
Löschen inkl.\ Balancieren & $O(\log n)$ \textbf{auch im Worst Case} \\
\end{tabular}
\par\vspace{3pt}
\textbf{Rekurrenz der Suche} (Praktikum 4): $T(n)=T(n/2)+c \Rightarrow O(\log n)$ — im Gegensatz
zum entarteten BST mit $T(n)=T(n{-}1)+c \Rightarrow O(n)$.
\end{defbox}

\begin{tipbox}{Klausurtipps}
\begin{itemize}
  \item \textbf{Der 8-P-Klassiker (Probeklausur):} Balancefaktor \emph{erklären} (Definition oben),
        Element \emph{einfügen} und \emph{Rebalancieren zeichnen}.
  \item Fürs volle Punktbild \textbf{explizit hinschreiben}: verletzten Knoten + seinen BF, das
        Kriterium ($BF$ des Kindes $\to$ einfach/doppelt), den Rotationstyp und die
        \textbf{End-Balancefaktoren} aller Knoten.
  \item \textbf{Rotationsrichtung} $=$ Gegenrichtung der Schieflage: linkslastig ($-2$) $\to$ nach
        \emph{rechts} rotieren.
  \item Nach jeder Rotation \textbf{Inorder-Check} — fängt Links/Rechts-Dreher sofort ab.
\end{itemize}
\end{tipbox}

\begin{pitbox}{Stolperfallen}
\begin{itemize}
  \item \textbf{Vorzeichen-Konvention:} $BF=$ \textbf{rechts $-$ links} (Folie 4) — viele Lehrbücher
        rechnen andersrum! Linkslastig $\Rightarrow$ \emph{negativ}.
  \item \textbf{Nach der Rechtsrotation wird der alte verletzte Knoten das \emph{rechte} Kind} des
        aufgestiegenen Knotens (er ist der größere von beiden) — nicht das linke!
  \item \textbf{Am untersten verletzten Knoten rotieren}, nicht an der Wurzel — auch wenn weiter
        oben ebenfalls $\pm 2$ steht.
  \item \textbf{Der mittlere Teilbaum ($b$) wechselt bei der Rotation die Seite} — beim Zeichnen
        gern vergessen.
  \item \textbf{Doppelrotation ist nicht „zweimal dieselbe Richtung``:} erst dreht sich das
        \emph{Kind} (begradigt den Knick), dann der verletzte Knoten in die Gegenrichtung.
\end{itemize}
\end{pitbox}

\clearpage
% ============================================================
%  KAP 08 — B-BÄUME
% ============================================================
\kapitel{08 — B-Bäume}{Ordnung, Suche, Einfügen mit Split, Löschen mit Verschieben/Verschmelzen}

\textbf{Kernidee.} BST/AVL nehmen an, dass alles in den Hauptspeicher passt. Bei großen
Datenmengen auf \textbf{externem Speicher} (Festplatte/SSD) zählt aber die Zahl der
\emph{Seitenzugriffe}: Ein B-Baum-Knoten $=$ eine \textbf{Seite} (page) $\Rightarrow$ dicke
Knoten mit vielen Schlüsseln $\Rightarrow$ flacher, breiter Baum $\Rightarrow$ wenige Zugriffe.
(\textbf{B} $=$ balanciert, breit, buschig, Bayer — \emph{nicht} binär; Bayer \& McCreight.)

\begin{defbox}{Ordnung $m$ — die Struktur-Regeln (Folien 5--7)}
\renewcommand{\arraystretch}{1.3}
\begin{tabular}{@{}p{4.6cm}p{4.2cm}p{4.2cm}@{}}
 & \textbf{Schlüssel} & \textbf{Kinder (innerer Knoten)} \\
\hline
je Knoten \textbf{mindestens} & $m-1$ & $m$ \\
je Knoten \textbf{höchstens} & $2m-1$ & $2m$ \\
\end{tabular}
\par\vspace{3pt}
\begin{itemize}
  \item \textbf{Immer: Kinder $=$ Schlüssel $+$ 1} — zwischen $i$ Schlüsseln liegen $i{+}1$ Verweise.
  \item \textbf{Wurzel-Ausnahme:} nur die Wurzel darf weniger als $m{-}1$ Schlüssel haben (min.\ 1).
  \item Schlüssel im Knoten sind \textbf{sortiert}; links von $e_j$ nur kleinere, rechts nur größere.
  \item \textbf{Alle Blätter liegen auf derselben Tiefe} — der Baum ist \emph{immer} balanciert.
  \item \textbf{Höhe:} $h \le \log_m \frac{n+1}{2}$ — wächst nur logarithmisch, Basis $m$.
\end{itemize}
\end{defbox}

\begin{defbox}{Suche (Folien 9--10)}
Kombination aus \emph{Suche im sortierten Array} (innerhalb des Knotens) und
\emph{Verweis-Verfolgen} (wie BST): im Knoten bis zum ersten Schlüssel $\ge key$ laufen —
gefunden $\to$ fertig; Blatt $\to$ NIL; sonst in den Unterbaum \emph{davor} absteigen.
\par\vspace{2pt}
\textbf{Effizienz:} externe Zugriffe $=$ besuchte Knoten $= O(\log_m n)$;
mit linearer Suche im Knoten insgesamt $O(\log_m n \cdot m)$.
\end{defbox}

\begin{defbox}{Rezept: Einfügen (Folien 11--17)}
\begin{enumerate}
  \item Neuer Schlüssel kommt \textbf{grundsätzlich in ein Blatt}.
  \item Beim Abstieg \textbf{jeden vollen Knoten ($2m{-}1$ Schlüssel) sofort splitten}:
        der \textbf{mittlere Schlüssel} wandert in den Elternknoten, links/rechts davon
        entstehen zwei Knoten mit je $m{-}1$ Schlüsseln.
  \item \emph{Warum beim Abstieg?} Der hochgereichte Schlüssel trifft so \textbf{nie auf einen
        vollen Elternknoten} (der wurde gerade geprüft) — ein einziger Abwärts-Durchlauf genügt.
  \item \textbf{Volle Wurzel} beim Start: erst splitten — neue Wurzel entsteht \emph{darüber},
        der Baum wächst \textbf{nach oben}. \emph{Das ist der einzige Weg, wie er höher wird.}
\end{enumerate}
\textbf{Effizienz:} $O(\log_m n)$ externe Zugriffe; Split selbst $O(1)$ Zugriffe / $O(m)$ intern.
\end{defbox}

\begin{defbox}{Rezept: Löschen (Folien 18--30)}
\textbf{Unterlauf} (underflow): Knoten fällt \emph{unter} $m{-}1$ Schlüssel $\to$ muss repariert
werden — ein Knoten verschwindet \emph{nie} einfach!
\begin{enumerate}
  \item Schlüssel im \textbf{Blatt}, kein Unterlauf $\to$ einfach entfernen.
  \item Blatt mit \textbf{Unterlauf}: kann ein \textbf{Geschwister abgeben} ($>m{-}1$ Schlüssel)
        $\to$ \textbf{Verschieben über den Elternknoten}: Trennschlüssel der Eltern rutscht
        \emph{runter}, äußerster Geschwister-Schlüssel rückt \emph{hoch}. Sonst
        $\to$ \textbf{Verschmelzen}: Knoten $+$ Trennschlüssel aus den Eltern $+$ Geschwister
        werden \emph{ein} Knoten (Eltern verlieren einen Schlüssel — kann sich nach oben fortsetzen).
  \item Schlüssel in einem \textbf{inneren Knoten}: der \textbf{Vorgänger} (größter Schlüssel im
        linken Kind) \textbf{rückt hoch} — geht direkt, wenn sein Knoten abgeben kann; sonst
        \textbf{Verschmelzen der beiden Nachfolger-Knoten} um die Lücke.
  \item Wird die \textbf{Wurzel leer} (letzter Schlüssel per Verschmelzen abgewandert):
        Wurzel entfernen $\to$ \textbf{Höhe sinkt um 1}.
\end{enumerate}
\end{defbox}

\begin{defbox}{Lösch-Trace (Ordnung $m=2$: min.\ 1, max.\ 3 Schlüssel) — alle vier Fälle}
\begin{Verbatim}[fontsize=\small,xleftmargin=6pt]
      [10 20]              loesche 12: kein Unterlauf -> einfach raus
     /   |   \
 [3 5] [12 15] [25 30]

      [ 5 20]              loesche 15: Unterlauf! Geschwister [3 5] kann
     /   |   \             abgeben -> VERSCHIEBEN ueber die Eltern:
   [3] [10]  [25 30]       10 runter ins Blatt, 5 hoch in die Wurzel

      [ 20 ]               loesche 3: Unterlauf, Geschwister [10] am Minimum
     /      \              -> VERSCHMELZEN: leer + Trenner 5 + [10] = [5 10]
  [5 10]   [25 30]         (Wurzel gibt die 5 ab -- Wurzel-Ausnahme ok)

      [ 10 ]               loesche 20 (innerer Knoten): VORGAENGER 10 aus
     /      \              dem linken Kind [5 10] rueckt hoch (kann abgeben,
   [5]     [25 30]         2 > Minimum 1)
\end{Verbatim}
\end{defbox}

\begin{defbox}{Praktikum 5 — Ordnung $m$ und Disk-I/O}
\begin{itemize}
  \item \textbf{Plattenzugriffe sinken} mit wachsendem $m$: Höhe $\log_m n$ schrumpft, weil die
        Log-Basis wächst (mehr Schlüssel pro geladener Seite).
  \item \textbf{Interne Vergleiche steigen} trotzdem: pro Knoten bis zu $m$ Vergleiche (lineare Suche)
        $\Rightarrow$ insgesamt $O(m\log_m n)$ pro Einfügung — der Faktor $m$ wächst schneller,
        als $\log_m n$ fällt. RAM-Vergleiche sind aber billig, Plattenzugriffe teuer $\to$ lohnt sich.
  \item \textbf{Optimales $m$ für 4-KB-Block} (Schlüssel und Kindzeiger je 8 Byte): Knoten voll
        $\Rightarrow$ $2m{-}1$ Schlüssel $+$ $2m$ Verweise:
        $8(2m{-}1) + 8 \cdot 2m = 32m - 8 \le 4096 \Rightarrow m \le 128{,}25 \Rightarrow \mathbf{m = 128}$.
\end{itemize}
\end{defbox}

\begin{tipbox}{Klausurtipps}
\begin{itemize}
  \item \textbf{Trace-Aufgaben:} Einfüge-Folge mit Splits bzw.\ Lösch-Folge zeichnen. Beim Löschen
        \emph{zuerst den Fall bestimmen}: Unterlauf? Kann ein Geschwister abgeben? Blatt oder innerer Knoten?
  \item \textbf{Selbstcheck nach jedem Schritt:} Kinder $=$ Schlüssel $+1$ in jedem Knoten?
        Alle Blätter auf gleicher Tiefe? Jeder Nicht-Wurzel-Knoten $\ge m{-}1$ Schlüssel?
  \item \textbf{Zahlen-Fragen:} Ordnung $m$ gegeben $\to$ min./max.\ Schlüssel \& Kinder nennen
        ($m{-}1$ / $2m{-}1$ bzw.\ $m$ / $2m$); Höhenformel $h\le\log_m\frac{n+1}{2}$ nachschlagen statt herleiten.
  \item \textbf{Merksatz:} \emph{„Unterlauf? Knoten bleibt! Erst beim Geschwister leihen (über die
        Eltern), sonst verschmelzen.``}
\end{itemize}
\end{tipbox}

\begin{pitbox}{Stolperfallen}
\begin{itemize}
  \item \textbf{Einen Unterlauf-Knoten nie einfach weglassen} — „Wurzel hat noch 2 Kinder`` ist das
        falsche Kriterium; prüfe \textbf{Kinder $=$ Schlüssel $+$ 1}. Ein 2-Schlüssel-Knoten \emph{muss}
        3 Kinder haben, sonst fehlt ein Wertebereich.
  \item \textbf{Verschieben läuft immer über den Elternknoten} — nie ein Schlüssel direkt vom
        Geschwister herüber (der Trennschlüssel würde die Sortierung verletzen).
  \item \textbf{Innerer Knoten: B-Baum-Folien nehmen den \emph{Vorgänger von links}} — beim BST
        (Kap.\ 06) war es der \emph{Nachfolger von rechts}. Nicht verwechseln!
  \item \textbf{Der B-Baum wächst an der Wurzel} (nur beim Wurzel-Split), nicht „wenn alle Knoten
        voll sind`` — und er schrumpft nur, wenn die Wurzel leer wird.
  \item \textbf{Beim Split wandert der mittlere Schlüssel hoch} — er bleibt \emph{nicht} zusätzlich
        in einem der beiden neuen Knoten.
\end{itemize}
\end{pitbox}

% ============================================================
%  Weitere Kapitel folgen — je Foliensatz \clearpage + \kapitel{}{}
% ============================================================

\end{document}
